\pagestyle{plain}

\chapter{CAPÍTULO 5: Auditoria Computacional e Aceitação Institucional de Dossiês de Salvaguarda de Saberes Agrícolas Tradicionais Quilombolas}


\begin{flushright}
Catuxe Varjão de Santana Oliveira  \footnote[1]{Programa de Pós-graduação em Ciências da Propriedade Intelectual-PPGPI/Universidade Federal de Sergipe-UFS. *E-mail:  \href{mailto:catuxe@academico.ufs.br}{\color{blue}{\underline{catuxe@academico.ufs.br}}}}
\\[0.3cm]
Luiz Diego Vidal Santos \footnote[2]{Universidade Estadual de Feira de Santana - UEFS, Brasil. E-mail: \href{mailto:ldvsantos@uefs.br}{\color{blue}{\underline{ldvsantos@uefs.br}}}}
\\[0.3cm]
Paulo Roberto Gagliardi \footnotemark[1]
\end{flushright}


% resumo em português — bookmark suprimido para manter hierarquia do capítulo
{\renewcommand{\PRIVATEbookmarkthis}[1]{}%
\begin{resumoumacoluna}
A invisibilidade documental de Saberes e Sistemas Agrícolas Tradicionais (SSAT) perante instâncias regulatórias configura uma falha de governança que perpetua a vulnerabilidade territorial das comunidades detentoras. Este estudo avalia a eficácia probatória e a aceitação institucional dos dossiês computacionais gerados pela arquitetura de salvaguarda desenvolvida nos capítulos precedentes, compreendendo fichas WOCAT-SLM-QBR adaptadas (Capítulo~3), perfis de vulnerabilidade meta-analíticos (Capítulo~2), escores ISB com classificação fuzzy de urgência (Capítulo~4) e metadados de rastreabilidade FAIR. A avaliação foi conduzida em duas etapas complementares. Na primeira, uma matriz de requisitos documentais foi construída cruzando os campos exigidos por quatro instâncias regulatórias (IPHAN, INPI, CGen e FAO-GIAHS) com os campos produzidos pelo sistema, gerando dossiês-piloto formatados para 3--5 práticas e saberes de comunidades quilombolas do Território de Identidade Semiárido Nordeste II (Bahia). Na segunda, um painel de 12--15 avaliadores com experiência comprovada em processos de salvaguarda patrimonial, registro de propriedade intelectual e acesso a patrimônio genético avaliou os dossiês mediante instrumento estruturado compreendendo dimensões de completude documental (razão campos presentes / campos requeridos) e aceitação institucional (escala Likert de 7 pontos). Adicionalmente, uma simulação de submissão regulatória foi conduzida para ao menos um saber documentado, identificando gargalos procedimentais não capturados pela avaliação em painel. Os indicadores primários compreendem taxa de completude documental ($\geq$ 85\%), mediana de aceitação institucional ($\geq$ 5,0/7), concordância interavaliadores ($W \geq$ 0,60) e conformidade FAIR ($\geq$ 80\%).

\vspace{\onelineskip}
\noindent
\textbf{Palavras-chave}: Auditoria computacional; Aceitação institucional; Dossiês de salvaguarda; Propriedade intelectual coletiva; Conformidade FAIR; Comunidades quilombolas.
\end{resumoumacoluna}}

\newpage

\section{Introdução}

Os capítulos precedentes desta tese desenvolveram um modelo integrado sequencial para decodificação, mensuração e priorização de Saberes e Sistemas Agrícolas Tradicionais (SSAT) quilombolas. O instrumento WOCAT-SLM-QBR adaptado (Capítulo~3) converteu saberes tácitos em variáveis linguísticas auditáveis. A meta-análise (Capítulo~2) sintetizou evidências sobre magnitudes de efeito das dimensões de vulnerabilidade biocultural. O Índice de Salvaguarda Biocultural (ISB), construído por integração TRI-fuzzy (Capítulo~4), classificou saberes segundo urgência de proteção em escala contínua 0--100. Contudo, a utilidade prática dessa cadeia computacional permanece condicionada à capacidade dos artefatos documentais por ela produzidos de atender aos requisitos probatórios das instâncias regulatórias responsáveis pela proteção formal do patrimônio imaterial e da propriedade intelectual coletiva.

O enquadramento regulatório brasileiro e internacional para proteção de saberes tradicionais envolve múltiplas instâncias com requisitos documentais heterogêneos e parcialmente sobrepostos. O Instituto do Patrimônio Histórico e Artístico Nacional (IPHAN) exige, para fins de Registro de Bens Culturais de Natureza Imaterial (Decreto 3.551/2000), dossiês contendo descrição do bem, território de ocorrência, comunidade detentora, modos de transmissão, ameaças e documentação audiovisual \cite{IPHAN2006}. O Instituto Nacional da Propriedade Industrial (INPI) demanda, para Indicações Geográficas (Lei 9.279/1996), delimitação territorial, caderno de normas, comprovação de notoriedade e vínculo geográfico \cite{Brasil1996}. O Conselho de Gestão do Patrimônio Genético (CGen), no âmbito da Lei 13.123/2015, requer prova de consentimento prévio informado, rastreabilidade do acesso ao patrimônio genético e conhecimento tradicional associado, e termos de repartição de benefícios \cite{Brasil2015}. O programa GIAHS da FAO exige plano de manejo dinâmico, evidência de agrobiodiversidade, comprovação de resiliência do sistema e governança comunitária ativa \cite{FAO2018}.

A questão que se coloca não é de valoração (quanto vale o saber), mas de governança documental, ou seja, se os dossiês produzidos pelo sistema computacional atendem aos requisitos probatórios dessas instâncias com completude, rastreabilidade e conformidade normativa suficientes para instruir processos de proteção formal. Nesse enquadramento, a aceitação institucional opera como proxy de eficácia do sistema de salvaguarda, sendo a hipótese testada que a documentação computacional auditável gerada pelo sistema (dossiês ISB, fichas WOCAT-SLM-QBR adaptadas e relatórios de inferência fuzzy) atende aos critérios de completude, rastreabilidade e conformidade normativa exigidos para instruir processos de salvaguarda junto a órgãos regulatórios, com taxa de aceitação institucional superior à de documentação convencional não padronizada.


\section{Materiais e Métodos}

O delineamento compreende quatro fases sequenciais, conforme detalhado na metodologia geral da tese.

\subsection{Construção da Matriz de Requisitos Regulatórios}

Foi construída uma matriz de requisitos documentais ($M_{r \times c}$) cruzando os $r$ campos exigidos por cada instância regulatória ($i \in \{IPHAN, INPI, CGen, GIAHS\}$) com os $c$ campos produzidos pelos módulos do sistema computacional de salvaguarda. Para cada instância, os requisitos formais de instrução processual foram mapeados a partir de consulta direta aos normativos vigentes e complementados por entrevistas semiestruturadas com 3--5 servidores ou consultores com experiência em instrução processual nessas instâncias.

O cruzamento produziu, para cada par instância-campo, uma classificação de cobertura $\zeta_{i,j} \in \{0, 0{,}5, 1\}$, correspondendo, respectivamente, a campo não coberto pelo sistema, campo parcialmente coberto (informação presente mas em formato ou granularidade insuficiente) e campo plenamente coberto. A taxa de cobertura sistêmica por instância foi calculada como $\bar{\zeta}_i = \frac{1}{r_i}\sum_{j=1}^{r_i}\zeta_{i,j}$.


\subsection{Geração de Dossiês-Piloto}

O sistema computacional de salvaguarda (módulos I a IV) foi aplicado a 3--5 práticas e saberes de comunidades quilombolas do Território de Identidade Semiárido Nordeste II (Bahia), gerando dossiês completos contendo perfil do saber (WOCAT-SLM-QBR adaptado), escore ISB com perfil de vulnerabilidade por dimensão (sistema fuzzy, Capítulo~4), perfil psicométrico TRI (Capítulo~4), metadados de rastreabilidade (proveniência, consentimento, cadeia de custódia documental) e material audiovisual georreferenciado.

Cada dossiê foi formatado em conformidade com os requisitos mapeados na fase anterior, produzindo versões específicas para cada instância regulatória (dossiê IPHAN, dossiê INPI, dossiê CGen, dossiê FAO-GIAHS). A conformidade FAIR (Findable, Accessible, Interoperable, Reusable) dos metadados de rastreabilidade foi avaliada mediante as métricas de \citeonline{Wilkinson2016}, com meta de $\geq$ 80\% de aderência.


\subsection{Painel de Avaliação Institucional}

Foi constituído um painel de 12--15 avaliadores com experiência comprovada em processos de salvaguarda, registro, indicação geográfica ou acesso a patrimônio genético, recrutados por conveniência e indicação em cadeia (\textit{snowball}) a partir de redes institucionais. O painel incluiu servidores do IPHAN (Coordenação de Registro), analistas do INPI (Diretoria de Contratos, Indicações Geográficas e Registros), consultores do CGen/MMA, e especialistas em sistemas agrícolas tradicionais com atuação em FAO ou organismos multilaterais.

Cada avaliador recebeu um dossiê-piloto formatado para a instância de sua especialidade e um instrumento estruturado de avaliação contendo duas dimensões. A primeira dimensão, Completude Documental ($C_d$), mensurou a presença e adequação de cada campo exigido pela instância regulatória, mediante escala dicotômica (presente/ausente) e escala Likert de adequação (1 a 5) para cada campo (Equação~\ref{eq:completude_cap5}).

\begin{equation}
C_d = \frac{\sum_{j=1}^{r_i} \mathbb{1}[\text{campo}_j \text{ presente}]}{r_i} \times 100\%
\label{eq:completude_cap5}
\end{equation}

A segunda dimensão, Aceitação Institucional ($A_i$), foi mensurada mediante escala Likert de 7 pontos (1, totalmente inadequado, a 7, plenamente adequado para instruir processo), avaliando clareza da evidência, rastreabilidade da documentação, conformidade normativa, suficiência probatória para deferimento e recomendação de uso do sistema. Foi incluída questão aberta para identificação de lacunas, redundâncias ou problemas de formato.


\subsection{Estudo de Caso Regulatório Simulado}

Para pelo menos um dos saberes documentados, foi conduzida uma simulação de submissão regulatória, preparando o dossiê completo e submetendo-o a avaliação por servidor ou consultor da instância-alvo em condições análogas às de tramitação real (sem protocolização formal). A simulação permitiu identificar gargalos procedimentais não capturados pela avaliação em painel, como exigências de formatação, prazos de complementação e interações entre campos do dossiê.


\subsection{Análise de Dados}

A taxa de completude documental foi expressa como percentual médio por instância regulatória ($\bar{C}_{IPHAN}$, $\bar{C}_{INPI}$, $\bar{C}_{CGen}$, $\bar{C}_{GIAHS}$), com intervalo de confiança de 95\%. A aceitação institucional foi analisada pela mediana e intervalo interquartil dos escores Likert, com teste de Friedman para comparação entre instâncias e análise de concordância interavaliadores mediante coeficiente de Kendall ($W$) (Equação~\ref{eq:kendall_cap5}).

\begin{equation}
W = \frac{12S}{k^2(n^3 - n)}
\label{eq:kendall_cap5}
\end{equation}

\noindent onde $S$ é a soma dos quadrados dos desvios dos totais de ranqueamento por instância, $k$ é o número de avaliadores e $n$ é o número de instâncias avaliadas.

A análise qualitativa das questões abertas seguiu codificação temática \cite{Braun2006} para identificação de categorias de inadequação e recomendações de ajuste, organizadas por instância regulatória. Os resultados foram triangulados com os achados da simulação regulatória para produzir um mapa de conformidade (\textit{compliance map}) sinalizando os campos do dossiê com maior e menor aderência aos requisitos institucionais.


\section{Resultados}

\textit{[Os resultados serão apresentados após a execução do protocolo de auditoria. A estrutura prevista segue abaixo.]}

\subsection{Matriz de Requisitos e Cobertura Sistêmica}

A matriz de requisitos será apresentada em formato tabular cruzando instâncias regulatórias e campos do sistema, com indicação visual de cobertura ($\zeta_{i,j}$). As taxas de cobertura sistêmica por instância ($\bar{\zeta}_i$) serão reportadas, identificando campos com cobertura parcial ou nula que demandem desenvolvimento adicional do sistema.

\subsection{Completude Documental dos Dossiês-Piloto}

As taxas de completude documental ($C_d$) serão reportadas por dossiê e por instância regulatória. A comparação entre instâncias permitirá identificar quais marcos regulatórios são melhor atendidos pelo modelo integrado e quais requerem módulos complementares. A meta de $\geq$ 85\% de campos presentes será avaliada formalmente.

\subsection{Aceitação Institucional}

As medianas e intervalos interquartis dos escores de aceitação institucional ($A_i$) serão reportados por dimensão de avaliação e por instância. O teste de Friedman identificará diferenças significativas entre instâncias. A concordância interavaliadores ($W$) indicará o grau de consenso do painel. A meta de mediana $\geq$ 5,0 será avaliada.

\subsection{Simulação Regulatória e Mapa de Conformidade}

Os resultados da simulação de submissão regulatória serão apresentados em formato narrativo, detalhando os gargalos procedimentais identificados. O mapa de conformidade consolidará achados quantitativos (completude, aceitação) e qualitativos (categorias temáticas das questões abertas, observações da simulação) em representação visual integrada.


\section{Discussão}

\textit{[A discussão será elaborada com base nos resultados obtidos. Os eixos analíticos previstos são apresentados abaixo.]}

As taxas de completude e aceitação serão interpretadas em termos de eficácia prática do sistema computacional como instrumento de governança documental, posicionando os resultados em relação a experiências internacionais de documentação de patrimônio imaterial e indicações geográficas. As lacunas identificadas na matriz de cobertura e no mapa de conformidade serão discutidas como oportunidades de desenvolvimento de módulos complementares do sistema.

A heterogeneidade de requisitos entre instâncias regulatórias será discutida como desafio estrutural para sistemas de documentação integrada, apontando para a necessidade de arquiteturas modulares e interoperáveis. As implicações para o design participativo (Capítulo~6) serão articuladas, demonstrando como os achados da auditoria informam o refinamento de interfaces e funcionalidades do sistema. As limitações, incluindo o tamanho do painel de avaliadores, a ausência de protocolização formal na simulação e a restrição ao marco regulatório brasileiro (com exceção do FAO-GIAHS), serão explicitadas.


\section{Conclusão}

\textit{[A conclusão será elaborada após os resultados. As entregas previstas são apresentadas abaixo.]}

O protocolo de auditoria computacional será apresentado como artefato de validação institucional do sistema de salvaguarda, demonstrando a capacidade do modelo integrado de produzir dossiês com eficácia probatória suficiente para instruir processos de proteção junto a instâncias regulatórias. As taxas de completude e aceitação serão posicionadas como indicadores de maturidade do sistema, informando a priorização de desenvolvimentos futuros e a aplicabilidade do protocolo como ferramenta de avaliação contínua.


%% ==============================================
%% NOTAS PARA PREENCHIMENTO APÓS EXECUÇÃO
%% ==============================================
%% [GAP CRÍTICO: Após execução do protocolo, inserir:
%%  - Tabela da matriz de requisitos regulatórios $M_{r \times c}$
%%  - Taxas de cobertura sistêmica $\bar{\zeta}_i$ por instância
%%  - Taxas de completude documental $C_d$ por dossiê e instância
%%  - Estatísticas descritivas de aceitação institucional ($A_i$)
%%  - Teste de Friedman e coeficiente $W$ de Kendall
%%  - Mapa de conformidade visual
%%  - Narrativa da simulação regulatória
%%  - Conformidade FAIR dos metadados]
